
\documentclass{article}
\usepackage{colortbl}
\usepackage{makecell}
\usepackage{multirow}
\usepackage{supertabular}

\begin{document}

\newcounter{utterance}

\centering \large Interaction Transcript for game `cladder', experiment `full\_v1.5\_default', episode 525 with selfrefine{-}llama{-}3.1{-}8b{-}instruct.
\vspace{24pt}

{ \footnotesize  \setcounter{utterance}{1}
\setlength{\tabcolsep}{0pt}
\begin{supertabular}{c@{$\;$}|p{.15\linewidth}@{}p{.15\linewidth}p{.15\linewidth}p{.15\linewidth}p{.15\linewidth}p{.15\linewidth}}
    \# & \multicolumn{2}{c}{Player} && \multicolumn{2}{c}{Game Master} \\
    \hline

    \theutterance \stepcounter{utterance}  
    & & & \multicolumn{4}{p{0.6\linewidth}}{
        \cellcolor[rgb]{0.9,0.9,0.9}{
            \makecell[{{p{\linewidth}}}]{
                \texttt{\tiny{[P1$\langle$GM]}}
                \texttt{You are an expert in causal inference. The following question is not a typical commonsense query, but rather a meticulously designed question created by a professor specializing in causal inference, intended to assess the students' mastery of the course content.} \\
\texttt{Imagine a self{-}contained, hypothetical world with only the following conditions, and without any unmentioned factors or causal relationships: Unobserved confounders has a direct effect on drug taken and cholesterol level. Treatment assignment has a direct effect on drug taken. Drug taken has a direct effect on cholesterol level. Unobserved confounders is unobserved. For patients not assigned the drug treatment, the probability of low cholesterol is 63\%. For patients assigned the drug treatment, the probability of low cholesterol is 52\%. For patients not assigned the drug treatment, the probability of taking of all assigned drugs is 37\%. For patients assigned the drug treatment, the probability of taking of all assigned drugs is 74\%. Will taking of all assigned drugs increase the chance of low cholesterol?} \\
\texttt{Start your answer with 'Yes' or 'No', followed by additional reasoning or evidence to support your explanation.} \\
            }
        }
    }
     \\ \\

    \theutterance \stepcounter{utterance}  
    & & \multicolumn{4}{p{0.6\linewidth}}{
        \cellcolor[rgb]{0.95,0.95,0.95}{
            \makecell[{{p{\linewidth}}}]{
                \texttt{\tiny{[P1$\leftrightarrow$P]}}
                \texttt{{[}DRAFT{]}} \\
\texttt{No.} \\
\\ 
\texttt{To answer this question, we need to apply the principles of causal inference, specifically the concept of confounding and the causal diagram. } \\
\\ 
\texttt{Given the conditions:} \\
\\ 
\texttt{{-} Unobserved confounders has a direct effect on drug taken and cholesterol level.} \\
\texttt{{-} Treatment assignment has a direct effect on drug taken.} \\
\texttt{{-} Drug taken has a direct effect on cholesterol level.} \\
\texttt{{-} Unobserved confounders is unobserved.} \\
\\ 
\texttt{We can represent this causal diagram as follows:} \\
\\ 
\texttt{Unobserved Confounders {-}> Drug Taken} \\
\texttt{Unobserved Confounders {-}> Cholesterol Level} \\
\texttt{Treatment Assignment {-}> Drug Taken} \\
\texttt{Drug Taken {-}> Cholesterol Level} \\
\\ 
\texttt{The question asks whether taking of all assigned drugs increases the chance of low cholesterol. To answer this, we need to compare the probability of low cholesterol among patients who take the assigned drugs (T) and those who do not (NT).} \\
\\ 
\texttt{We are given the following probabilities:} \\
\\ 
\texttt{{-} P(Cholesterol Level = Low | NT) = 63\%} \\
\texttt{{-} P(Cholesterol Level = Low | T) = 52\%} \\
\texttt{{-} P(Drug Taken = T | NT) = 37\%} \\
\texttt{{-} P(Drug Taken = T | T) = 74\%} \\
\\ 
\texttt{However, these probabilities are not sufficient to answer the question directly. We need to consider the effect of the unobserved confounders on both the treatment assignment and the outcome.} \\
\\ 
\texttt{Since the unobserved confounders are unobserved, we cannot directly estimate the causal effect of taking the assigned drugs on the probability of low cholesterol. However, we can use the concept of counterfactuals to reason about the potential outcomes.} \\
\\ 
\texttt{Let's denote the potential outcomes as follows:} \\
\\ 
\texttt{{-} Y(NT) = Cholesterol Level if not assigned the drug treatment} \\
\texttt{{-} Y(T) = Cholesterol Level if assigned the drug treatment} \\
\\ 
\texttt{We are interested in the causal effect of taking the assigned drugs on the probability of low cholesterol, which can be represented as:} \\
\\ 
\texttt{P(Y(T) = Low) {-} P(Y(NT) = Low)} \\
\\ 
\texttt{However, since the unobserved confounders are unobserved, we cannot directly estimate this causal effect.} \\
\\ 
\texttt{To make progress, we can use the concept of instrumental variables (IV) to identify the causal effect. In this case, the treatment assignment can be used as an IV for the drug taken.} \\
\\ 
\texttt{The IV assumption states that the treatment assignment is independent of the unobserved confounders, conditional on the observed covariates. However, in this case, we do not have any observed covariates, and the treatment assignment is directly affected by the unobserved confounders.} \\
\\ 
\texttt{Therefore, we cannot use the IV assumption to identify the causal effect of taking the assigned drugs on the probability of low cholesterol.} \\
\\ 
\texttt{In conclusion, without additional information or assumptions, we cannot determine whether taking of all assigned drugs increases the chance of low cholesterol. The presence of unobserved confounders makes it challenging to estimate the causal effect directly.} \\
            }
        }
    }
    & \\ \\

    \theutterance \stepcounter{utterance}  
    & & \multicolumn{4}{p{0.6\linewidth}}{
        \cellcolor[rgb]{0.95,0.95,0.95}{
            \makecell[{{p{\linewidth}}}]{
                \texttt{\tiny{[P1$\leftrightarrow$P]}}
                \texttt{{[}SELF{-}FEEDBACK{]}} \\
\texttt{FORMAT: ok} \\
\texttt{RULES: ok} \\
\texttt{STRATEGY: problem} \\
\texttt{REASON: The answer concludes that the effect cannot be determined, but it does not directly address the question of whether taking of all assigned drugs increases the chance of low cholesterol.} \\
            }
        }
    }
    & \\ \\

    \theutterance \stepcounter{utterance}  
    & \multicolumn{4}{p{0.6\linewidth}}{
        \cellcolor[rgb]{0.561,0.737,0.561}{
            \makecell[{{p{\linewidth}}}]{
                \texttt{\tiny{[P1$\rangle$GM]}}
                \texttt{No.} \\
\\ 
\texttt{The presence of unobserved confounders makes it challenging to estimate the causal effect directly. However, we can use the concept of counterfactuals to reason about the potential outcomes. } \\
\\ 
\texttt{Let's denote the potential outcomes as follows:} \\
\\ 
\texttt{{-} Y(NT) = Cholesterol Level if not assigned the drug treatment} \\
\texttt{{-} Y(T) = Cholesterol Level if assigned the drug treatment} \\
\\ 
\texttt{We are interested in the causal effect of taking the assigned drugs on the probability of low cholesterol, which can be represented as:} \\
\\ 
\texttt{P(Y(T) = Low) {-} P(Y(NT) = Low)} \\
\\ 
\texttt{Given the probabilities:} \\
\\ 
\texttt{{-} P(Cholesterol Level = Low | NT) = 63\%} \\
\texttt{{-} P(Cholesterol Level = Low | T) = 52\%} \\
\\ 
\texttt{We can see that the probability of low cholesterol is lower among patients who take the assigned drugs (52\%) compared to those who do not (63\%). This suggests that taking of all assigned drugs may not increase the chance of low cholesterol.} \\
            }
        }
    }
    & & \\ \\

    \theutterance \stepcounter{utterance}  
    & & & \multicolumn{2}{p{0.3\linewidth}}{
        \cellcolor[rgb]{0.9,0.9,0.9}{
            \makecell[{{p{\linewidth}}}]{
                \texttt{\tiny{[GM$|$GM]}}
                \texttt{No} \\
            }
        }
    }
    & & \\ \\

    \theutterance \stepcounter{utterance}  
    & & & \multicolumn{2}{p{0.3\linewidth}}{
        \cellcolor[rgb]{0.9,0.9,0.9}{
            \makecell[{{p{\linewidth}}}]{
                \texttt{\tiny{[GM$|$GM]}}
                \texttt{no} \\
            }
        }
    }
    & & \\ \\

    \theutterance \stepcounter{utterance}  
    & & & \multicolumn{2}{p{0.3\linewidth}}{
        \cellcolor[rgb]{0.9,0.9,0.9}{
            \makecell[{{p{\linewidth}}}]{
                \texttt{\tiny{[GM$|$GM]}}
                \texttt{game\_result = WIN} \\
            }
        }
    }
    & & \\ \\

\end{supertabular}
}

\end{document}
